\section{Limitations}\label{sec:limitations}

\textbf{Training wall-clock.} CliffordSTF training lags the fastest spherical-harmonic baselines (MACE, EquiformerV2) by roughly two-to-threefold per step at matched parameter budget, with the autograd graph through the dual-track scaffold dominating the overhead (\S\ref{sec:results_scalar}, Supplementary~\ref{supp:wallclock}). Further headroom of an estimated $40$--$60\%$ is plausible through full \texttt{torch.compile} engagement (currently blocked upstream by \texttt{torch\_scatter} custom-op tracing under dynamic shapes) and fused Triton kernels for the STF$_2$/STF$_3$ products. The STF$_2$-only variant (Table~\ref{tab:ablation}) may suffice for long MD trajectories at lower per-step cost. At $10^7$ parameters, per-step cost is roughly threefold the $10^6$-parameter regime, in line with the additional capacity.

\textbf{Training stability on scalar benchmarks.} QM9 and Molecule3D required gradient clipping to prevent early-training divergence from the bilinear cross-track coupling. A more principled cross-track initialisation or dedicated STF normalisation would likely address this without the clipping.

\textbf{Fixed-budget protocol.} Targeting cross-architecture comparability rather than per-architecture tuning means several baselines (notably MACE, ICTP, and GotenNet) underconverge relative to their published recipes (per-model reasoning in Supplementary~\ref{supp:results_caveats}). Our catalysis results characterise the sub-10M-parameter regime rather than the leaderboard frontier, where published OC20/OC22 numbers use $10$--$100\times$ more parameters.

\textbf{Absolute directional fidelity gap.} While CliffordSTF improves rMD17 directional fidelity by an order of magnitude over base Clifford, absolute fidelity remains below fully spherical-harmonic architectures: EquiformerV2 at $L\!\geq\!2$ attains $\sim\!0.96$ aggregate force-cosine similarity, versus $0.551$ for CliffordSTF. Closing this remaining gap without spherical harmonics, Wigner-$D$ matrices, or CG tables is a natural extension.

\textbf{Scope of the directional analysis.} Our diagnostic is restricted to rMD17 small organics in vacuum. Whether the $L\!\leq\!1$ clustering at poor cosine alignment persists on periodic systems, larger molecules, or solvated condensed phases remains open, though CliffordSTF's catalysis results (\S\ref{sec:results_catalysis}) suggest the phenomenon generalises beyond the rMD17 setting.

\section{Discussion}\label{sec:discussion}

This paper joins two equivariant lineages that have developed largely in parallel. Geometric algebra treats the multivector as the primary substrate and derives rotational covariance from the geometric product, while irreducible Cartesian tensors treat symmetric traceless tensors as the primary substrate and derive covariance from algebraic symmetry. Our results suggest neither substrate alone is sufficient for atomistic force prediction at the per-edge bilinear level: the Clifford geometric product reaches $L \leq 1$ exactly, Cartesian symmetric traceless tensors cover $L \geq 2$ but forfeit the natural scalar and pseudoscalar structure of the Clifford algebra, and only the hybrid scaffold (\S\ref{sec:results_scalar}, Table~\ref{tab:ablation}) attains the order-of-magnitude directional gain that motivates the architecture. The broader implication is that representational choice in equivariant architectures should be governed by the angular content one actually wants to represent, not by loyalty to a particular algebraic formalism; Theorem~\ref{thm:cg_completeness} makes this concrete by showing that the geometric product and closed-form STF operations together span the same $\mathrm{SO}(3)$-irrep content as spherical harmonics through $L{=}3$. Within the hybrid, the ablation further isolates the $L \geq 2$ scaffold, rather than any single coupling flag, as the dominant accuracy driver, and identifies cross-track coupling as a training-stability lever; this is consistent with the theoretical framing, since the scaffold is what unlocks representational capacity while the coupling governs whether the gradients flowing through it behave predictably. We see promising extensions along three axes: higher $L_{\max}$ through iterated STF constructions, fused hardware kernels to close the wall-clock gap against spherical-harmonic baselines, and full $\mathrm{O}(3)$ equivariance through explicit parity tracking of polar and axial vector channels (Remark~\ref{rem:so3_vs_o3}). We release the full PyTorch implementation of CliffordSTF along with all baseline methods and benchmark hyperparameters at \texttt{https://github.com/KurbanIntelligenceLab/CliffordSTF}
